\documentclass[10pt,a4paper]{amsart}
\usepackage[margin=19mm]{geometry}
\usepackage{amsmath,amssymb,mathtools}
\usepackage[scr=rsfs,cal=euler]{mathalfa}
\usepackage{xfrac}
\usepackage[unicode,hidelinks,bookmarks=false]{hyperref}
\newcommand{\ie}{i.e.\ }
\newcommand{\cf}{cf.\ }
\newcommand{\resp}{respectively\ }
\newcommand{\Subsection}{\subsection}
\providecommand{\nobreakdash}{\nobreak\hbox{-}}
\setlength{\emergencystretch}{2em}
\multlinegap0pt
\newtheorem{theorem}{Theorem}
\title{The $\theta$ invariant recovers the Rozansky-Overbay invariant}
\author{Ramana Murugesan}
\date{Source: arXiv:2604.27029v1 (CC BY 4.0)}
\begin{document}
\maketitle

\noindent\emph{Source and licence.} The $\theta$ invariant recovers the Rozansky-Overbay invariant. Version arXiv:2604.27029v1 (CC BY 4.0), distributed by arXiv under the Creative Commons Attribution 4.0 International licence, http://creativecommons.org/licenses/by/4.0/, as recorded in the arXiv licence metadata for this exact version and shown on the abstract page https://arxiv.org/abs/2604.27029v1. Full licence notice: \texttt{licenses/arxiv-2604.27029-CC-BY-4.0.txt}.\par\smallskip

\noindent\textit{Editorial scope.} Counted unit: the article's theorem thm:main (source0.tex lines 29--31). The article's complete original proof of it is delivered with it (source0.tex lines 97--127, one \texttt{proof} environment of 31 lines, which the article places away from the statement). No mathematics is written, repaired or re-derived: the statement and the proof are the article's own lines, in the article's own notation, and no retained line is substituted or re-flowed.  Retained uncounted as the source-local context the proof needs: lines 49--51.  Nothing was omitted from any retained span, so the package carries no omission record.  No retained line contains a citation, so the wrapper carries no bibliography and no bibliography entry.  No retained line contains a figure environment, an image-inclusion command or drawing code, so no image is delivered and the package declares no asset dependency.  Editorial formatting only, in addition: an AMS \texttt{amsart} preamble that restates the article's own declarations and macros used by the retained lines, namely the environments \texttt{theorem} on separate counters, together with the standard packages those lines need.  The retained environments print in this document as Theorem 1 in order of appearance, whereas the article prints the same items with the numbering of its own full document.  The complete native source of the article as distributed in the arXiv e-print for this exact version is bundled unchanged as \texttt{sources/199440/source0.tex}.
\begin{equation}
    A=I-\sum_{c\in X} (T^s E_{i,i+1}+(1-T^s)E_{i,j+1}+E_{j,j+1}) \label{A}
\end{equation}

\begin{theorem}\label{thm:main}
    Let $K$ be a knot. Then $\rho_1(K)=-\theta(K)|_{T_1\leftarrow 1,T_2\leftarrow T}$
\end{theorem}

\begin{proof}[Proof of Theorem \ref{thm:main}]
    Let $D$ be a upright knot diagram representing $K$.
    Since $T_3=T_1T_2=T$ we have
\begin{align*}g_{ii}&=g_{2ii}=g_{3ii}\\
\Delta&=\Delta_2=\Delta_3\\
\Delta_1&=\Delta|_{T\leftarrow1}=1
\end{align*}
where the last equality comes from the normalization of $\Delta$. Since $F_2\equiv0$
and $F_3(k)=\varphi_k(g_{kk}-1/2)$ 
it suffices to show that $R_1(c)=-F_1(c)$ for each $c=(s,i,j)$.
With the above equations,
\begin{equation*}
    F_1(c)=\frac{s}{2}[1+2(-1+T^s)g^2_{ji}+2g_{ji}(1+g_{ij}-2g_{jj})+2g_{ii}(-1-(-1+T^s)g_{ji}+g_{jj})] \label{eq:f1}
\end{equation*}

To show that this equals $-R_1(c)$, we need to eliminate the
$g_{j,j+1}$ and $g_{j+1,j}$ terms in $R_1(c)$.
Recall that $G=A^{-1}$. Multiplying \eqref{A} by $G$,
\begin{align*}
1&=(AG)_{jj}=g_{jj}-g_{j+1,j}\\
0&=(GA)_{j,j+1}=g_{j,j+1}-(1-T^s)g_{ji}-g_{jj}\end{align*}
Thus:
\begin{align*}
    R_1(c)&=s\left[\vphantom{\frac12}g_{ji}((g_{jj}-1)+(g_{jj}+(1-T^s)g_{ji})-g_{ij})\right.\\
    &\quad-g_{ii}(g_{jj}
    \left.+(1-T^s)g_{ji}-1)-\frac{1}{2}\right]\\
        &= \frac{s}{2}[-1+2(1-T^s)g_{ji}^2+2g_{ii}(-g_{jj}-(1-T^s)g_{ji}+1)\\
        &\quad+2g_{ji}(2g_{jj}-1-g_{ij})]\\
        &= -F_1(c)
\end{align*}
\end{proof}

\end{document}
